\documentclass{amsart}
\usepackage{amsmath,amssymb,amsthm,hyperref}

\theoremstyle{definition}
\newtheorem{definition}{Definition}
\newtheorem{thesis}{Thesis}
\theoremstyle{plain}
\newtheorem{claim}{Claim}
\newtheorem{lemma}{Lemma}
\newtheorem{proposition}{Proposition}
\newtheorem{objection}{Objection}
\theoremstyle{remark}
\newtheorem{remark}{Remark}

\begin{document}

\title{Scientific Objectivity and the Separation of the Descriptive and
Justificatory Questions about Induction}
\maketitle

\begin{abstract}
We give a precise account of ``scientific objectivity'' that keeps the
\emph{descriptive} question about inductive reasoning (quid facti?) and the
\emph{normative} question (quid juris?) strictly distinct. The account
identifies objectivity not with the acceptance of an inductive procedure, nor
with the (Humean) demand for an unconditional vindication of induction, but
with a \emph{relational, framework-explicit} property: a procedure is objective
relative to an explicitly adopted, publicly inspectable inductive framework
when its appraisals are determined by that framework's content alone and the
framework's adoption is itself open to non-parochial, agent-invariant
assessment. We are careful \emph{not} to claim a solution to Hume's problem;
our thesis is conditional, and we prove that on \emph{any} non-trivial choice
of background framework the descriptive set and the justificatory set come
apart, that the justificatory question is preserved rather than dismissed, and
that classification and justification are never conflated. We then defend the
distinction against the strongest forms of deflationism --- synchronic
community-relativism \emph{and} Peircean limit-of-inquiry pragmatism --- and
close the dialectic by answering the deflationist's return objections.
\end{abstract}

\section*{Methodological preface: what is and is not claimed}

This problem is conceptual, not a theorem of a fixed formal theory; ``proof''
here means an explicit, gap-free argument from stated definitions. To keep the
argument honest we flag two temptations and refuse them.

\emph{First}, we do \emph{not} claim to solve Hume's problem of induction.
Hume's problem is that any premise strong enough to vindicate induction (e.g.\
a uniformity-of-nature postulate) is itself either an inductive claim (regress)
or a stipulation (only conditional warrant). We therefore do not assert that
the set of objective procedures is non-empty \emph{absolutely}. Instead we make
all justification \emph{relative to an explicitly adopted framework}, and we
prove the separation theses for every admissible framework. This converts the
Humean obstacle from a hidden gap into an explicit parameter.

\emph{Second}, we do not assume a unique correct ``logic of partial support.''
Which confirmation logic is correct (a Carnapian $\lambda$-continuum, a
subjective-Bayesian setup, Popperian corroboration, $\dots$) is itself a
contested normative question. We therefore index everything to a chosen
framework and prove that our claims hold uniformly across the admissible range,
so that nothing rests on a privileged choice.

\section{Formalisation of the Problem}

\begin{definition}[Inductive framework]\label{def:framework}
An \emph{inductive framework} is a triple $\Phi=(L,\mathcal{M},\vdash_\Phi)$ in
which
\begin{itemize}
\item $L$ is a fixed background logic of full consequence (e.g.\ classical
first-order logic) together with, where applicable, a fixed calculus of partial
support (e.g.\ the Kolmogorov probability axioms);
\item $\mathcal{M}$ is a finite, explicitly stated list of \emph{methodological
postulates} (a uniformity assumption, a simplicity ordering, an admissible
class of priors, a projectibility constraint, $\dots$), each formulated as a
proposition;
\item $\vdash_\Phi$ is the relation ``$\Phi$ certifies $\,\cdot\,$'' defined as
provability in $L$ from $\mathcal{M}$.
\end{itemize}
A framework is \emph{non-trivial} if $\mathcal{M}$ does not contain, or entail
under $L$, the proposition ``a procedure is warranted iff it is endorsed by the
prevailing community.'' (We exclude only this single degenerate framework,
which builds deflationism in by fiat.)
\end{definition}

\begin{remark}[Why frameworks, and why this is not circular]\label{rem:circ}
Making $\Phi$ explicit blocks two charges at once. (i)~\emph{Circularity}: the
predicate ``objective relative to $\Phi$'' defined below quantifies over agents
who reason by $L$, but $L$ is a \emph{parameter} fixed before the
descriptive\slash justificatory contrast is drawn, exactly as ``valid
argument'' is defined relative to a fixed semantics without circularity. (ii)~
\emph{Concealed Humean gap}: by listing $\mathcal{M}$ openly we make visible
that the warrant we certify is \emph{conditional} on $\mathcal{M}$; we never
pretend to have derived $\mathcal{M}$ from nothing. All theorems below are
proved ``for every non-trivial $\Phi$,'' so no privileged $L$ or $\mathcal{M}$
is presupposed.
\end{remark}

\begin{definition}[Inductive procedure]\label{def:procedure}
An \emph{inductive procedure} is a rule $\Pi$ mapping a body of evidence $E$
(a finite set of observation propositions) and a hypothesis $h$ to an appraisal
$\Pi(h\mid E)$ in a fixed appraisal space $A$. We do \emph{not} require $\Pi$
to be ``purely syntactic.'' On the contrary, a procedure is individuated by the
descriptive vocabulary in which it frames evidence and hypotheses; in
particular two procedures may agree on all logical form yet differ in which
predicates they treat as projectible. This is essential: by Goodman's lesson
no purely syntactic fact distinguishes a projectible predicate (``green'')
from a non-projectible one (``grue''), so projectibility must be located in the
framework's postulates $\mathcal{M}$, not in syntax.
\end{definition}

\begin{definition}[The descriptive question, quid facti?]\label{def:facti}
For which procedures $\Pi$ do human inquirers, as a matter of empirical fact,
employ $\Pi$ (or an extensional approximation) with apparent practical success?
Write
\[
\mathcal{F}=\{\Pi:\Pi\text{ is, as a matter of fact, employed by successful
inquirers}\}.
\]
Membership in $\mathcal{F}$ is contingent and settled by observation of
behaviour. ``Success'' here is read with a \emph{content-bearing} criterion $S$
(predictive accuracy, instrumental reliability, $\dots$), not as a synonym for
``endorsed''; see Lemma~\ref{lem:noncoincide}.
\end{definition}

\begin{definition}[Conditional justification, quid juris?]\label{def:juris}
Relative to a framework $\Phi$, a \emph{justification} of $\Pi$ is a finite
$L$-proof from premises drawn from $\mathcal{M}$ (and from logical truths)
whose conclusion is ``$\Pi$ is warranted relative to the cognitive aim fixed by
$\Phi$.'' Write
\[
\mathcal{J}_\Phi=\{\Pi:\Phi\vdash_\Phi\text{``}\Pi\text{ is warranted''}\}.
\]
Membership in $\mathcal{J}_\Phi$ is a logical matter relative to $\Phi$,
settled by checking $L$-derivability, never by observation of behaviour.
\emph{We do not claim $\mathcal{J}_\Phi$ is non-empty for an unconditionally
justified $\Phi$; the warrant it records is explicitly conditional on
$\mathcal{M}$.}
\end{definition}

\section{The Proposed Account of Objectivity}

\begin{definition}[Agent-invariance]\label{def:invariance}
A justification $J$ (an $L$-proof from $\mathcal{M}$) is \emph{agent-invariant}
iff every premise of $J$ lies in $\mathcal{M}\cup\{L\text{-truths}\}$ and no
inference step adverts to the identity, location, history, or
community-membership of the reasoner. Equivalently: $J$ is checkable by any
agent who has accepted $\Phi$, with the same verdict, because $L$-derivability
from a fixed premise set is agent-independent.
\end{definition}

\begin{definition}[Public inspectability of the framework]\label{def:public}
A framework $\Phi$ is \emph{publicly inspectable} iff $\mathcal{M}$ is stated
explicitly (no tacit appeal to ``what we do''), so that any agent can (a)~see
which postulates are in force and (b)~contest any postulate by offering
$L$-reasons against it. Public inspectability is what makes the adoption of
$\Phi$ itself answerable to non-parochial assessment, rather than a brute fiat.
\end{definition}

\begin{definition}[Scientific objectivity --- the proposed account]
\label{def:objectivity}
A procedure $\Pi$ is \emph{objective relative to $\Phi$}, written
$\Pi\in\mathcal{O}_\Phi$, iff
\begin{enumerate}
\item[(O1)] $\Phi$ is publicly inspectable and non-trivial
(Definitions~\ref{def:framework},~\ref{def:public});
\item[(O2)] there is an agent-invariant justification of $\Pi$ in $\Phi$
(so $\Pi\in\mathcal{J}_\Phi$); and
\item[(O3)] $\Pi$'s appraisals are fixed by the content of $E,h$ and
$\mathcal{M}$ alone --- i.e.\ $\Pi(h\mid E)$ does not vary with the identity of
the user.
\end{enumerate}
\emph{Objectivity simpliciter} is then the relational fact ``objective relative
to some publicly inspectable non-trivial framework whose postulates are not
themselves rationally indicted within $\Phi$.'' We make no stronger,
framework-free claim.
\end{definition}

The central move: objectivity is a \emph{logical property of a warrant relative
to an openly avowed framework} (does an agent-invariant $\Phi$-proof exist?),
\emph{not} a statistical property of the population of inquirers, and \emph{not}
the unconditional Humean vindication. The framework parameter is precisely
where the Humean debt is paid openly rather than hidden.

\section{Adequacy: Conditions (a)--(c)}

Throughout, ``accepted practice'' is membership in $\mathcal{F}$. Fix an
arbitrary non-trivial, publicly inspectable framework $\Phi$.

\begin{proposition}[Condition (a): no reduction to, and no definability from,
description]\label{prop:a}
For every such $\Phi$: \emph{(i)} the membership criterion of
$\mathcal{O}_\Phi$ contains no behavioural predicate, hence $\mathcal{O}_\Phi$
is not definable from $\mathcal{F}$ by substitution of its definiens; and
\emph{(ii)} $\mathcal{O}_\Phi\neq\mathcal{F}$ in extension.
\end{proposition}

\begin{proof}
\emph{(i) Non-definability.} The defining conditions (O1)--(O3) quantify only
over: the syntactic form of $\mathcal{M}$ and of $L$-proofs, the content of
$E,h$, and the absence of agent-indexicals. None of these is a behavioural
predicate; ``$\Pi$ is employed by successful inquirers'' never occurs in the
definiens. Therefore there is no substitution instance turning the definiens of
$\mathcal{O}_\Phi$ into a formula all of whose non-logical vocabulary is drawn
from the behavioural vocabulary used to define $\mathcal{F}$. (This is the exact
sense of ``not definable from $F$'' demanded by condition (a): the criterion of
objectivity does not \emph{consist in} a fact about practice.) Non-definability
is thus established directly from the form of the definition, independently of
any extensional facts.

\emph{(ii) Extensional difference}, with explicit instances.
\begin{itemize}
\item \emph{$\mathcal{F}\setminus\mathcal{O}_\Phi$.} Let $\Pi_{\text{par}}$ be a
procedure whose appraisals depend on the user's community --- e.g.\ a procedure
that outputs ``confirmed'' precisely when the in-group's elders have so
declared. This violates (O3) outright (its output varies with the user), so
$\Pi_{\text{par}}\notin\mathcal{O}_\Phi$; yet a community may use it and
prosper, so $\Pi_{\text{par}}\in\mathcal{F}$. (Note this instance does
\emph{not} rely on Goodman's grue, precisely because grue and green are
syntactically on a par; the parochialism here is in the explicit dependence on
who applies $\Pi$, which (O3) targets directly.)
\item \emph{$\mathcal{O}_\Phi\setminus\mathcal{F}$.} Take $\Pi^\ast$ to be the
appraisal ``output the $\Phi$-probability of $h$ given $E$ to $10^{6}$ decimal
places.'' For suitable $\Phi$ (one whose $\mathcal{M}$ fixes a computable prior)
$\Pi^\ast$ satisfies (O1)--(O3), so $\Pi^\ast\in\mathcal{O}_\Phi$; yet no actual
inquirer carries out that computation, so $\Pi^\ast\notin\mathcal{F}$.
\item \emph{$\mathcal{O}_\Phi\cap\mathcal{F}$} (a procedure both used and
$\Phi$-warranted, e.g.\ ordinary $\Phi$-conditionalisation on a tractable case)
and the joint complement are likewise instantiable.
\end{itemize}
All four cells are non-empty, so $\mathcal{O}_\Phi\neq\mathcal{F}$ and neither
inclusion is definitional. With (i) this establishes (a).
\end{proof}

\begin{proposition}[Condition (b): the justificatory question is preserved,
not postponed]\label{prop:b}
For every non-trivial publicly inspectable $\Phi$, the account treats the quid
juris? question as genuine and answerable, and does not discharge it by mere
postponement.
\end{proposition}

\begin{proof}
By (O2), $\Pi\in\mathcal{O}_\Phi$ requires $\Pi\in\mathcal{J}_\Phi$, so the
notion is unintelligible unless ``is $\Pi$ $\Phi$-warranted?'' is a real
question; indeed $\mathcal{O}_\Phi\subseteq\mathcal{J}_\Phi$. It might be
objected (cf.\ the critic's G1) that conditionalising warrant on $\mathcal{M}$
merely \emph{postpones} quid juris? to ``is $\mathcal{M}$ warranted?''. We meet
this head-on. The account does not \emph{suppress} that further question; by
public inspectability (Definition~\ref{def:public}) each postulate in
$\mathcal{M}$ is exposed for exactly the second-order justificatory scrutiny
quid juris? demands. What the account refuses --- correctly, given Hume --- is
to \emph{pretend} that scrutiny terminates in an unconditional vindication. The
honest content of quid juris? is: \emph{given that some framework must be
adopted to reason inductively at all, which frameworks survive agent-invariant,
publicly inspectable criticism of their postulates?} That question is preserved
in full force as the criterion for ``objectivity simpliciter'' in
Definition~\ref{def:objectivity}. Thus quid juris? is neither dismissed
(condition (b)) nor falsely declared solved (the preface's honesty
constraint).
\end{proof}

\begin{proposition}[Condition (c): no conflation of classification and
justification]\label{prop:c}
For every non-trivial publicly inspectable $\Phi$, the factual classification
(membership in $\mathcal{F}$) and the justificatory question (membership in
$\mathcal{J}_\Phi$ or $\mathcal{O}_\Phi$) are settled by disjoint evidential
bases and never inter-substituted.
\end{proposition}

\begin{proof}
Deciding $\Pi\in\mathcal{F}$ requires observation of behaviour and of the
content-bearing success criterion $S$ (Definition~\ref{def:facti}); this datum
never enters (O1)--(O3). Deciding $\Pi\in\mathcal{J}_\Phi$ (resp.\
$\mathcal{O}_\Phi$) requires checking $L$-derivability from $\mathcal{M}$
(resp.\ also (O1),(O3)); this never appeals to behaviour. The two procedures
take disjoint inputs (a behavioural record vs.\ a proof object) and yield
verdicts by disjoint criteria. By Proposition~\ref{prop:a}(ii) the outputs
diverge in all four cells, so no inference ``$\Pi\in\mathcal{F}$ therefore
$\Pi\in\mathcal{O}_\Phi$,'' or its converse, is licensed. Objectivity is still
secured (non-vacuously, relative to $\Phi$) because $\mathcal{O}_\Phi$ is
populated whenever a procedure carries an agent-invariant $\Phi$-proof and is
user-independent --- a condition checkable by inspecting proofs, with no
head-counting. This establishes (c).
\end{proof}

\begin{thesis}[Conditional adequacy]\label{thesis:adequacy}
For every non-trivial publicly inspectable framework $\Phi$,
Definition~\ref{def:objectivity} yields a notion of scientific objectivity
satisfying (a), (b), and (c); and it does so without claiming a framework-free
solution to Hume's problem.
\end{thesis}

\begin{proof}
Immediate from Propositions~\ref{prop:a}--\ref{prop:c}, the universal
quantification over $\Phi$ neutralising the ``which logic of support?'' worry,
and the explicit conditionalisation of warrant on $\mathcal{M}$ neutralising
the ``empty/question-begging'' worry.
\end{proof}

\section{Defence against the Deflationary Alternative}

\begin{objection}[Deflationism about the distinction]\label{obj:deflation}
The descriptive/justificatory distinction is misplaced: to call $\Pi$ justified
is just to report endorsement by successful practice, so $\mathcal{F}$,
$\mathcal{J}_\Phi$ and $\mathcal{O}_\Phi$ coincide and
Definition~\ref{def:objectivity} mislabels a sociological fact as a logical
property. Two strengths must be addressed: \emph{(D-syn)} synchronic
community-relativism (warrant $=$ current endorsement), and \emph{(D-lim)}
Peircean limit pragmatism (warrant $=$ endorsement at the ideal limit of
inquiry).
\end{objection}

\begin{lemma}[Non-coincidence under any content-bearing success criterion]
\label{lem:noncoincide}
If ``success'' carries content $S$ beyond ``is endorsed,'' then
$\mathcal{F}\neq\mathcal{J}_\Phi$ is possible, so $\mathcal{F}=\mathcal{J}_\Phi$
is a substantive, not analytic, identity --- and the distinction is therefore a
genuine question.
\end{lemma}

\begin{proof}
The deflationist needs ``successful practice'' lest $D$ endorse every fad, so
$S$ is invoked. Either (i) $S$ reduces to ``endorsed,'' making the qualifier
vacuous (handled by Lemma~\ref{lem:selfapp}); or (ii) $S$ has
endorsement-independent content. In case (ii) a community can endorse a
$\Pi$ that fails $S$ (lucky streaks, biased self-audit) and fail to endorse one
that meets $S$; and $\Phi$-warrant tracks $L$-derivability from $\mathcal{M}$,
not the endorsement record. Hence the extensions can diverge, so identity is
substantive.
\end{proof}

\begin{lemma}[Self-application dilemma, with the question-begging reply blocked]
\label{lem:selfapp}
The deflationist thesis $D$ cannot be advanced without either lacking rational
force or presupposing the agent-invariant warrant it denies; and the natural
``you beg the question'' reply does not rescue it.
\end{lemma}

\begin{proof}
Ask whether $D$ is warranted.

\emph{Case 1 ($D$ warranted in the deflationist's own sense).} Then $D$'s
correctness is a contingent report of one community's dispositions; it supplies
no reason binding on an interlocutor who does not share them. $D$ then has the
status of a datum for $\mathcal{F}$, not a refutation in $\mathcal{J}_\Phi$.

\emph{Case 2 ($D$ warranted simpliciter, i.e.\ something the interlocutor ought
to accept on reasons).} Then asserting $D$ presupposes a notion of ``ought to
accept on reasons'' not constituted by endorsement --- i.e.\
agent-invariant warrant (Definition~\ref{def:invariance}) --- which $D$ denies;
$D$ is self-refuting.

\emph{Blocking the question-begging reply.} The deflationist may protest that
Case~1 begs the question by presupposing community-transcendent ``reasons.''
But the reply collapses into the dilemma. If \emph{no} norm is
endorsement-independent, then the very dispute is, by the deflationist's lights,
a clash of dispositions adjudicable by no reason --- which is exactly the loss
of rational force Case~1 diagnoses, now embraced. The non-deflationist needs no
``view from nowhere''; it suffices that \emph{some} norm (modus ponens; the
probability axioms; ``do not affirm $p$ and $\neg p$'') be correct
independently of being endorsed. A deflationist who grants even one such norm
has granted agent-invariance; one who grants none cannot present $D$ \emph{as
correct} rather than \emph{as ours}, and so cannot indict our account.

\emph{Blocking the quietist reading.} A Rortyan quietist may say ``ought to
accept'' just means ``warranted by our shared practice,'' dissolving Case~2.
But then $D$ asserted to the non-deflationist is, again, merely a report of the
quietist's practice (Case~1): it gives the non-deflationist, whose practice
differs, no reason to revise. The quietist re-reading does not add a third
horn; it relocates Case~2 into Case~1. The dichotomy is therefore exhaustive
for present purposes: every construal of ``$D$ is warranted'' is either
force-bearing-and-self-refuting or force-free.
\end{proof}

\begin{lemma}[Indispensability for criticism: against \emph{both} D-syn and
D-lim]\label{lem:criticism}
A standard of warrant not constituted by endorsement is required to represent
ordinary scientific self-criticism; this defeats synchronic deflationism, and
limit-pragmatism does not evade it.
\end{lemma}

\begin{proof}
\emph{Against D-syn.} Under ``warrant $=$ current endorsement,'' ``currently
endorsed-and-successful practice is in error'' is contradictory. But the
history of science is replete with endorsed, locally successful procedures
later judged unwarranted on reflective grounds (ad hoc rescues,
suppressed-evidence inferences). Representing such immanent self-correction
requires an endorsement-independent standard --- supplied here by
agent-invariant $\Phi$-justification (Definition~\ref{def:invariance}): a
community can ask whether its own endorsed $\Pi$ carries an agent-invariant
$\Phi$-proof and discover it does not.

\emph{Against D-lim.} The Peircean replaces current endorsement with
endorsement at the ideal limit of inquiry. This concedes our central point in
disguised form. First, ``the limit of inquiry'' is not an observable
sociological fact about any actual community, so D-lim is \emph{not} a
descriptive reduction of warrant to practice; it already abandons the
deflationist's quid facti?/quid juris? collapse. Second, to give ``limit''
content the Peircean must specify which inquiry-policies count as truth-tracking
(else the limit is undefined or trivial); that specification is precisely a set
of methodological postulates $\mathcal{M}$ assessed for truth-conduciveness ---
i.e.\ a framework $\Phi$ evaluated by agent-invariant criteria. Thus D-lim,
made precise, \emph{is} a special case of our account (objectivity relative to
the framework whose postulates encode the limit-conditions), not a rival to it.
Either way an endorsement-independent standard is in play, so the distinction
is indispensable.
\end{proof}

\begin{lemma}[Closing the dialectic: the deflationist's return objections]
\label{lem:return}
The deflationist's strongest rejoinders --- ``your agent-invariant warrant is a
myth because (R1) your $\mathcal{M}$ faces Hume's regress, and (R2) there is no
unique logic of support'' --- do not re-establish the collapse of the
distinction.
\end{lemma}

\begin{proof}
\emph{R1 (regress).} We have not claimed unconditional warrant; warrant is
explicitly conditional on $\mathcal{M}$ (Definition~\ref{def:juris}, preface).
The regress shows only that $\mathcal{O}_\Phi$ is relative to $\Phi$, which we
assert, not that ``objective'' \emph{means} ``endorsed.'' Crucially, a
\emph{conditional} but agent-invariant warrant is still not a sociological fact:
$\Phi\vdash_\Phi\Pi$ is true or false independently of who endorses $\Phi$. So
R1, even if granted entirely, leaves $\mathcal{O}_\Phi\neq\mathcal{F}$
(Proposition~\ref{prop:a}) intact and hence does not collapse the distinction.

\emph{R2 (no unique support logic).} Our theorems are universally quantified
over admissible $\Phi$ (Thesis~\ref{thesis:adequacy}); they do not require a
unique $L$. If several support logics are admissible, we obtain a \emph{family}
of objectivity predicates $\{\mathcal{O}_\Phi\}$, each separating from
$\mathcal{F}$. Plurality of frameworks is not endorsement-relativism: the
choice among admissible $\Phi$ is itself constrained by agent-invariant,
publicly inspectable criticism of postulates (Definition~\ref{def:public}), not
by head-counting. Hence R2 yields framework-pluralism about objectivity, which
is still inconsistent with the deflationist's identification of warrant with
endorsement.

Since R1 and R2 are the deflationist's only routes back from
Lemmas~\ref{lem:noncoincide}--\ref{lem:criticism}, and both fail, the dialectic
is closed.
\end{proof}

\begin{thesis}[Refutation of the alternative]\label{thesis:refute}
The view that the descriptive/justificatory distinction is misplaced is
untenable, and the framework-explicit account of
Definition~\ref{def:objectivity} is to be preferred.
\end{thesis}

\begin{proof}
By Lemma~\ref{lem:noncoincide} the identity $\mathcal{F}=\mathcal{J}_\Phi$ is
substantive, not analytic, so the distinction is a genuine question. By
Lemma~\ref{lem:selfapp} (with the question-begging and quietist replies
blocked) $D$ either lacks rational force or self-refutes. By
Lemma~\ref{lem:criticism} an endorsement-independent standard is indispensable
for scientific self-criticism, against synchronic \emph{and} limit-pragmatist
deflationism. By Lemma~\ref{lem:return} the deflationist's return objections
(Humean regress, plurality of support logics) do not re-collapse the
distinction. Hence the objection fails on every count, while by
Thesis~\ref{thesis:adequacy} the proposed account meets (a)--(c) for every
admissible framework. It is therefore adequate and preferable.
\end{proof}

\section{Summary}

Scientific objectivity is, on this account, a \emph{relational, framework-
explicit} property: a procedure $\Pi$ is objective relative to a publicly
inspectable, non-trivial inductive framework $\Phi$ when $\Pi$ has an
agent-invariant $\Phi$-justification and user-independent appraisals
(Definition~\ref{def:objectivity}), and objective simpliciter when it is so
relative to some framework whose postulates survive agent-invariant, publicly
inspectable criticism. We deliberately do \emph{not} solve Hume's problem:
warrant is conditional on the avowed postulates $\mathcal{M}$, and our claims
are proved uniformly over all admissible frameworks, so nothing rests on a
unique logic of support. On this basis: (a)~objectivity does not reduce to or
get defined from description --- its criterion contains no behavioural
predicate and the extensions diverge in all four cells
(Proposition~\ref{prop:a}); (b)~the justificatory question is preserved in full,
indeed sharpened to ``which frameworks survive non-parochial criticism?'', and
is not falsely declared solved (Proposition~\ref{prop:b}); (c)~classification
and justification use disjoint evidential bases and are never inter-substituted
(Proposition~\ref{prop:c}). Deflationism --- synchronic and limit-pragmatist
alike --- is refuted by the non-analyticity of $\mathcal{F}=\mathcal{J}_\Phi$, a
self-application dilemma immune to the question-begging and quietist replies,
the indispensability of an endorsement-independent standard for self-criticism,
and the failure of the deflationist's Humean and pluralist return objections
(Lemmas~\ref{lem:noncoincide}--\ref{lem:return}).
\end{document}
